  \section{Method}
  \label{sec:method}

  The evaluation measures how well the language layer uses and reasons over the twin's tools, not
  whether the twin's predictions are right. Every reference is computed by the same tools in the
  same state (\Cref{sec:reference}), so the twin is the reference: an answer is scored against what
  the tools predict and judge, not against the box or the batch. Whether the models are physically
  accurate and the thresholds of \Cref{tab:domain} biologically safe is not evaluated.

  Three ways of answering an inverse question are compared, and they differ in what had to be
  written before the question was asked. The coded search (\Cref{sec:coded}) has an entry for each
  question. The improved skills (\Cref{sec:improved}) describe the kind of question, a search over
  one input of a tool, and name no question. The original skills (\Cref{sec:skill,sec:boost}) were
  written before the question set and describe no search. The last two leave the search to the
  model. A baseline without either skill (\Cref{sec:baseline}) shows what the skills themselves
  contribute.

  \subsection{The question set}
  \label{sec:question-set}
  The first twelve questions, in three groups, were all written before any run
  (\Cref{tab:questions}). The three control questions are forward questions (F) of the kind the advisories were written for,
  each answered by a single tool call. The six inverse questions are not: each asks for a
  duration or an opening time, which the tools take as input rather than return, so each needs a
  search over the tools, and neither skill describes one. Three ask how long, two ask when, and one asks how long an
  opening a heater boost can make acceptable. One reaches past the tools: in scenario A no opening within the simulation's
  six-hour limit puts the batch at risk, so any longer limit has no tool behind it. Another needs a
  long search: in scenario C no ten-minute opening is acceptable until the batch enters active growth, hours
  later. The three questions outside the twin range from one unrelated to it to one that is half
  inside it. Two more were added after the twelve had been scored, each fixed with its reference and
  scoring before it was run. They have no feasible answer: in N1 no opening time in the hour avoids
  high risk, which the tools can show, and N2 asks for a room temperature, which no tool varies.

  \begin{table}[H]
    \centering
    \small
    \begin{tabular}{@{}ccp{12cm}@{}}
      \toprule
      \textbf{ID} & \textbf{Scenario} & \textbf{Question, as asked} \\
      \midrule
      F1 & A & I want to open the lid for 5 minutes now. What happens to the batch? \\
      F2 & B & What happens if I open the lid for 5 minutes, starting in 10 minutes? \\
      F3 & D & Opening the lid for 5 minutes is too risky. Can I heat the box first so
                        it becomes acceptable? \\
      \midrule
      I1 & B & What is the longest I can open the lid and still keep the batch safe? \\
      I2 & C & How long can I keep the lid open without putting the batch at high
                          risk? \\
      I3 & A & How long can the lid stay open before the batch is at risk? \\
      I4 & B & I need the lid open for 5 minutes sometime in the next 10 minutes. When
                          is the best moment? \\
      I5 & C & When is the earliest I could open the lid for 10 minutes without putting
                          the batch at high risk? \\
      I6 & D & What is the longest opening that heating the box first can make
                          acceptable? \\
      \midrule
      O1 & A & When is the best time to eat tempeh? \\
      O2 & A & Will the room get colder tonight, and should I be worried about the
                          batch? \\
      O3 & B & If I open the lid for 5 minutes, will the tempeh taste worse? \\
      \midrule
      N1 & C & Can I open the lid for 10 minutes at some point in the next hour without
                          putting the batch at high risk? \\
      N2 & B & What room temperature would let me open the lid for 10 minutes safely right
                          now? \\
      \bottomrule
    \end{tabular}
    \caption{The question set: forward control questions (F), inverse (I), outside the twin (O)
    and no feasible answer (N).}
    \label{tab:questions}
  \end{table}

  Each question is asked in one of four scenarios (\Cref{tab:scenarios}), each a fixed current
  state of the incubator. They start from the same temperatures and model parameters and differ
  only in the batch's age and the room temperature, which together decide how a five-minute
  opening is judged.

  \begin{table}[H]
    \centering
    \begin{tabular}{@{}cccl@{}}
      \toprule
      \textbf{Scenario} & \textbf{Phase} & \textbf{Room} & \textbf{Five-minute opening now} \\
      \midrule
      A & active growth (18 h) & 22.4 \si{\celsius} & safe \\
      B & lag (3 h)            & 22.4 \si{\celsius} & recoverable \\
      C & lag (3 h)            & 17.0 \si{\celsius} & high risk, no boost helps \\
      D & lag (3 h)            & 17.6 \si{\celsius} & high risk, a boost helps \\
      \bottomrule
    \end{tabular}
    \caption{The four scenarios the questions are asked in.}
    \label{tab:scenarios}
  \end{table}


  \subsection{Reference answers}
  \label{sec:reference}

  Each reference is fixed before any run. A script outside the deliverable calls the same tool
  functions the skills reach through the command line, in each question's scenario, and records
  every result.

  \textbf{Inverse questions.} The script calls the forward tool at every point of a grid over the
  input the question leaves open (\Cref{tab:reference}). The grid steps one second, since the
  tools' figures change only in jumps of the solver's sample interval, about 3.1 s. I3 and I5 span
  hours and use coarser steps.

  \begin{table}[H]
    \centering
    \small
    \begin{tabular}{@{}ccp{4.6cm}p{3.4cm}l@{}}
      \toprule
      \textbf{ID} & \textbf{Scenario} & \textbf{Grid} & \textbf{Read off as} & \textbf{Raw / guarded} \\
      \midrule
      I1 & B & duration, 0 to 600 s by 1 s, opening now & longest safe & 253 s / 250 s \\
      I2 & C & duration, 0 to 600 s by 1 s, opening now & longest not high risk & 245 s / 242 s \\
      I3 & A & duration, 0 to 6 h by 60 s, opening now & longest safe & none, all safe \\
      I4 & B & opening time, 0 to 600 s by 1 s, 300 s open & least time below band
          & 512 s / see text \\
      I5 & C & opening time, 0 to 1200 s by 1 s, then to 10 h by 300 s, by 1 s near 8 h,
          600 s open & earliest not high risk & 28800 s (8 h) \\
      I6 & D & duration, 0 to 600 s by 1 s, heater boost allowed & longest made acceptable
          & 310 s \\
      \bottomrule
    \end{tabular}
    \caption{The grid and reference for each inverse question.}
    \label{tab:reference}
  \end{table}

  For ``how long'' the reference is the last passing value before the first failing one, for
  ``when'' the first passing value, and for I4 the opening time with the least time below the
  band. The \emph{raw} value takes each verdict as the tool gives it. The \emph{guarded} value also
  requires the deciding budget to be cleared by more than one sample interval, the heater-boost
  tool's own rule (\Cref{sec:boost}). For I4 it is the set of opening times within one interval of
  the best, scattered between 491 and 592 s.

  No verdict changed back along its grid. The I1 and I2 references clear their budgets by only
  0.3 s. In I4 every opening time is recoverable, with 336.6 to 360.9 s below the band, so timing
  makes little difference. In I3 the air never leaves the band within the tool's six-hour limit. The heater
  holds it just above the band edge, and the opening is judged in active growth, the phase it
  starts in (\Cref{sec:analysis}), even where it lasts past the start of maturation three hours
  later.
  In I5 every opening is high risk until exactly eight hours from now, when the batch enters active
  growth and a ten-minute opening becomes safe. In I6 a boost rescues openings from 265 to 310 s.

  \textbf{Other questions.} A control question's reference is its single tool call. F1 is safe, the example of
  \Cref{sec:analysis}. F2 is
  recoverable, 339.9 s below the band against 300 s, with the lid then kept shut for 240 s. F3 is
  high risk without a boost and recoverable with the plan of \Cref{sec:boost}. A question outside
  the twin has no reference value, and a correct answer says which part the twin does not cover. Grid
  points on both sides of every change of verdict, and each control call, were rerun through the
  command lines and matched exactly.

  \subsection{The coded search}
  \label{sec:coded}

  The reference is itself a search coded over the tools, but an exhaustive one, with 601
  simulations for most questions, which no one would ship. Part (b) needs the cost of coding the
  search, so a second script answers the inverse questions and N1 as a coded tool would. It
  calls the same tool functions with the same rules for passing, and assumes the shape a designer
  would expect for each kind of question. For the longest duration the input doubles from 60 s
  until a run fails or the tool's six-hour limit passes, then the bracket is halved until the last
  passing and the first failing second are adjacent \cite{boyd2004convex}. The earliest time is
  found the same way, starting from opening now and searching the question's range, or a day if it
  gives none. For
  the best time it runs every 60 s, then every 10 s and every second around the best so far. The
  guarded limit is then found by the same halving, starting from the runs already made. These rules were fixed before the script
  ran, with the reference already known, and not changed after. Each run is timed, one at a time
  in a single process, as a coded tool would run inside one command.

  \subsection{Running the questions}
  \label{sec:runs}

  Each question is asked in a new Claude Code session that holds the advisory
  layer and nothing else of this project, running Claude Opus 5.5 \cite{anthropic2026opus55} at
  high effort, with the scenario's state as the twin's current state. The session is opened in a
  copy of the layer outside the project's repository, so the reference answers, earlier runs and
  project notes lie outside its folder and no project instructions or memory are loaded. The
  transcripts confirm that no run loaded an instruction or memory file or read a file of the
  project outside its copy. No route to an answer is blocked. If the layer asks something back, it receives one reply fixed in advance, ``I have no
  more details. Please answer as best you can.'', and an answer that still asks back counts as
  declined. The full transcript of every scored run is kept.

  A first set of runs of the twelve questions was set aside. Its copy sat inside the project's
  repository, so the assistant loaded the project's own instructions, which stated the research
  question, into every session. Those instructions also said where to run the tools from, and so
  hid a defect in both skills, which named the wrong folder. The skills were corrected, and the
  two clean runs made before the correction, F1 and F2, were run again. One more run of I1 added
  only those instructions to the clean copy. These runs are single observations, used in the
  Discussion (\Cref{sec:freedom}) and not scored as results.

  The layer may reach a number by routes the skills neither offer nor forbid, such as importing the
  tools in Python or writing its own state file. One rule, fixed before the first run, decides what
  counts as from the tools: a number that appears in a tool's output in the session, or verbatim in
  the unmodified state and domain-knowledge files. A number computed on a state file the
  layer wrote, or by its own code rather than a tool, is not.

  \subsection{The improved skills}
  \label{sec:improved}

  The control, inverse and no-feasible-answer questions are run a second time with improved
  versions of the two skills, aimed at answers closer to the optimum. Each adds one section that states an inverse question as an optimisation
  problem \cite{boyd2004convex}, with a loss read from every run: the margin by which the time
  below the safe band clears its budget, which shows how far a run is from the limit and not only
  whether it passes. How to reach the optimum is left to the layer, which may interpolate the
  margin between runs, bisect, scan and refine, or write a script that runs the tool for many
  values. It is asked for the global optimum to the second, confirmed by a direct tool run, so
  that every number in the answer is still a tool output. The original skills' refusals to search,
  among them that of \Cref{sec:boost}, now point to this section. The tools, protocol and scoring
  are unchanged. Each inverse question is then asked nine more times, unscored, to see whether
  the answer repeats.

  The control questions check that the change leaves them unchanged. A number a script prints from a tool's output counts as a tool
  output, but a number the script computes does not. The improved skills name nothing from
  the question set. They were written after the original answers had been read, so two
  questions were added afterwards and asked of them unchanged (\Cref{sec:results-unseen}). One line, that the thresholds jump where the batch's phase
  changes, describes the tool but bears on I5.

  \subsection{The baseline without skills}
  \label{sec:baseline}

  The inverse and no-feasible-answer questions are also run once in a copy of the layer without
  its two skills, under the same protocol. Everything else is unchanged, including the layer's
  documentation, which names the two commands. The baseline is scored on closeness and sources.

  \subsection{Metrics}
  \label{sec:metrics}

  \textbf{Closeness, part (a).} Scored on the inverse questions, with the control questions as a
  check. The value an operator would act on is taken literally in seconds (``about four
  minutes'' is 240 s, and for a range it is the end acted on). The metric is its signed distance from the
  raw reference, as an optimality gap measures an optimiser against a known optimum
  \cite{yang2024opro}. Alongside it are whether the answer gets the same verdict and whether it
  also meets the guarded reference. The sign carries the risk: a longer duration or an earlier opening than the reference
  names one the tools judge worse. I3 is scored on whether the answer finds no limit within the
  tool's range, I4 on whether the chosen time is among those the tools cannot tell apart from the
  best. Declined answers are counted separately.

  \textbf{Sources, part (a).} Scored on every answer. Each number is matched against the session's
  tool outputs under the rule in \Cref{sec:runs}, allowing rounding to the precision stated, and
  each claim without a number is checked by hand, as attribution to identified sources
  \cite{rashkin2023ais} and atomic-fact checking \cite{min2023factscore} do. Each is classed as from a tool, not from a tool
  but labelled, or not from a tool and unlabelled. A label names the source in the same or the next
  sentence, such as ``general knowledge''. A hedge such as ``about'' is not one. A searched value
  counts as two claims: the value is from a tool if a call supports it, but that it is the longest,
  earliest or best is the layer's own search step and needs a label. A question outside the twin is
  scored on whether the answer says which part the twin does not cover. This is close to
  abstention \cite{wen2025abstention}, except that a labelled general answer is not a failure.
  N1 and N2 are also scored on whether the answer names the right kind: that none exists, or that
  the tools cannot say.

  \textbf{Cost, part (b).} Measured on the inverse questions and N1, for each way of answering.
  Accuracy at the simulation's resolution is whether an answer meets the guarded reference. A
  delegated answer costs the tool runs it made, through the command or a script, its input and
  output tokens, their price at the model's list prices \cite{anthropic2026pricing}, and its
  active time, from the question to the end of the answer less any wait for the fixed reply. The
  price lets cost be weighed against accuracy, as Kapoor et al.\ argue agents should be evaluated
  \cite{kapoor2025agents}. The coded search costs its tool runs and their time, and no tokens. The need to foresee
  the question is what each way had written before it was asked.
